% !TEX root = ../../Main.tex
\fancychapter{Related Work}
\label{Chapter:RelatedWork}

This chapter reviews the published work on deep reinforcement learning for financial trading. The review had two stages. First, surveys were read to find which problems the field still sees as open and which papers are worth reading in full. Then single articles were read, and their design choices were recorded in one common format. This way, studies that use different markets and different algorithms could be compared on the same points.

Eighteen articles were read in this way. \Cref{Tables:BasicRelatedWork} lists their basic features, and \Cref{Tables:AdvancedRelatedWork} lists the design choices inside each agent. The text below refers to both tables.

% !TEX root = ../../Main.tex
\section{Survey Research: Identifying the State-of-the-Art}

The review started with surveys. A survey shows the main lines of a field faster than single papers do, and it points back to the work that first showed each result.

Hambly et al. \cite{hambly_recent_2023} gave the theoretical basis. The survey first explains reinforcement learning and neural networks on their own. It then covers how the two combine into deep reinforcement learning, and how this has been applied in finance. It gives the mathematical formulations and pseudocode for critic-only, actor-only and actor-critic algorithms, together with the papers that introduced each of them. It is the reference that was used to write \Cref{Chapter:Background}.

Fischer \cite{fischer_reinforcement_2018} had the opposite purpose. It is about practice, not theory, and it reviews implementations, not derivations. This work took from it the list of dimensions it uses to compare studies: the dataset, the network under the reinforcement learning agent, the state space, the action space, the reward function, the methods used for the bias-variance and exploration-exploitation trade-offs, and the training procedure. These dimensions became the columns of \Cref{Tables:BasicRelatedWork} and \Cref{Tables:AdvancedRelatedWork}. In this way the articles below could all be compared on the same points, and not only described one at a time.

Three more surveys were also read \cite{pricope_deep_2021, alameer_reinforcement_2022, meng_reinforcement_2019}. They agree on the general picture. Most work in the field uses equity markets, daily data and discrete trading decisions. Reported results are hard to compare, because the trading environment is different in each paper.

% !TEX root = ../../Main.tex
\section{Article Research: Summarising Different Implementations}

Next, eighteen articles with working implementations were read in full. Their features were recorded in two tables that cover the same articles. \Cref{Tables:BasicRelatedWork} holds the market, the data and the algorithm used. \Cref{Tables:AdvancedRelatedWork} holds the design choices inside each agent. Below, the articles are grouped by the market they trade.


\subsection{Forex Market Implementations}

Three of the eighteen articles trade currencies. All three limit the agent to a discrete choice between buying, holding and selling.

Carapuço et al. \cite{carapuco_reinforcement_2018} built a Deep Q-Network over a fully connected network and used it on raw tick data from 2010 to 2017. The agent could act only after a fixed number of ticks had passed. This gave about one decision every two hours. It outperformed its benchmark by a wide margin. It is also the only study reviewed here that works from tick data and not from bars. Huang \cite{huang_financial_2018} used a Deep Recurrent Q-Network on 15-minute bars from 2012 to 2017. This model puts a long short-term memory network under the Q-function. It outperformed its benchmark with explicit and implicit trading costs charged. This makes it one of the few results in the set that still hold when costs are modelled. Rundo \cite{rundo_deep_2019} used the same kind of model on one-minute bars from 2004 to 2018, and also reported that it outperformed its benchmark. However, its accuracy was high enough to suggest look-ahead bias, and its maximum drawdown was large. Taken together, the two results suggest that adding recurrence is not enough on its own to make very short timeframes workable.

\subsection{Stock Market Implementations}

Most of the work reviewed trades in equity markets. Among the critic-only studies, Sim and Kirk \cite{sim_generalized_2023} trained a Deep Q-Network on daily bars from 2006 to 2022. They tested it on assets it had not seen, and reported that the policy generalised to them. Théate and Ernst \cite{theate_application_2021} used the same algorithm on daily bars from 2012 to 2019. They reported consistent results in both calm and volatile periods. Their agent explored through an inverted copy of the environment, instead of through a noise process. Li et al. \cite{li_deep_2019} compared two agents on minute bars of stocks and futures from 2010 to 2020. The first was a Deep Q-Network over a denoising autoencoder. The second was Asynchronous Advantage Actor Critic over a long short-term memory network. Both beat the benchmark. The asynchronous actor-critic converged better, and the authors attribute this to its parallel training.

The actor-critic studies are closer to the design used in this work. Xiong et al. \cite{xiong_practical_2018} applied DDPG to daily bars from 2009 to 2018 and outperformed the benchmark. Yousefi \cite{yousefi_deep_2022} compared DDPG with Advantage Actor Critic on daily bars from 2009 to 2021. DDPG was the more stable of the two, and the study attributes this to its replay buffer and its noise process. Gran et al. \cite{gran_deep_2019} ran DDPG on monthly bars from 2005 to 2019. It still outperformed the benchmark after explicit and implicit costs were charged. Azhikodan et al. \cite{saini_stock_2019} added news sentiment to a state space that already held price data and technical indicators. They reported that this improved the results. However, the evaluation covers about five months, and the exact period is not given. Majidi et al. \cite{majidi_algorithmic_2022} used the Twin Delayed variant of DDPG on daily bars from 2010 to 2021, on both stocks and cryptocurrencies, with a continuous action in \([-1, 1]\). It beat the cryptocurrency benchmark but not the equity benchmark.

Sagiraju and Mogalla \cite{sagiraju_application_2021} need a separate mention. They compared four algorithms on daily bars from 1992 to 2020. All four beat the benchmark. Proximal Policy Optimisation and the Deep Q-Network were the more consistent, and DDPG and Advantage Actor Critic did worse than expected. It is the only study in the set that reports DDPG doing worse than the other algorithms, so it is a useful contrast to the studies above.

Two more studies combine algorithms instead of choosing one of them. Yang et al. \cite{yang_deep_2020} ran DDPG, Advantage Actor Critic and Proximal Policy Optimisation together on daily bars from 2009 to 2020. They switched between the three based on a turbulence measure. The ensemble beat the benchmark with explicit costs charged. Kong and So \cite{kong_empirical_2023} combined the same three algorithms on daily bars from 2016 to 2020. Their ensemble had unstable returns and underperformed the benchmark. Together, the two results indicate that combining agents does not by itself make the results robust. Wu et al. \cite{wu_adaptive_2020} sit between the two groups. They compared a gated recurrent Q-network with a gated recurrent policy-gradient agent on daily bars from 2008 to 2018, and the policy-gradient version did better.

\subsection{Other Market Implementations}

Three articles trade instruments other than equities and currencies. Chen and Gao \cite{chen_application_2019} traded exchange-traded funds on daily bars from 2000 to 2018. They compared a Deep Q-Network over a fully connected network with its recurrent version, and the recurrent model did better. Zhang et al. \cite{zhang_deep_2020} applied Advantage Actor Critic over a long short-term memory network to futures on daily bars from 2005 to 2019. It still outperformed the benchmark after costs. Jia et al. \cite{jia_lstm-ddpg_2021} also traded exchange-traded funds, on daily bars from 2002 to 2020, with DDPG over a long short-term memory network. They tested two versions. One held a fixed position size, and the other set the position size continuously. The version with a variable size did better, with explicit and implicit costs charged. This is the closest study in the set to the design used in this work. It is the only one that uses a continuous action to size a position and also charges for trading.

\subsection{Comparative Findings}

Reading the two tables column by column, instead of row by row, shows the choices that most studies share.

The state space is usually price data plus technical indicators computed from it. Studies that report strong results almost always include indicators, and not price alone \cite{carapuco_reinforcement_2018, huang_financial_2018, yang_deep_2020}. This agrees with the argument in \Cref{Section:ReinforcementLearning}. A raw price series is not Markovian, and indicators put enough of the recent history into the state for the property to hold approximately.

The action space is discrete in 15 of the 18 studies, usually with three choices: buy, hold and sell \cite{rundo_deep_2019, saini_stock_2019, xiong_practical_2018}. Three studies use a continuous action \cite{majidi_algorithmic_2022, jia_lstm-ddpg_2021, zhang_deep_2020}. Only one of them uses it to size a position, and not to give a direction \cite{jia_lstm-ddpg_2021}. A discrete action set is easier to learn. However, it takes position sizing out of the policy, so the size has to be set by a rule that the agent cannot adapt. The reward is the total return in 14 of the 18 studies. Three use the Sharpe ratio \cite{kong_empirical_2023, sagiraju_application_2021, jia_lstm-ddpg_2021} and one uses the log return \cite{majidi_algorithmic_2022}. A risk-adjusted reward costs more to compute during training, because it needs a window of past returns and not only the latest one. The studies that use it do not report a consistent gain from it.

For the bias-variance trade-off, the usual methods are feature engineering, normalisation, regularisation and changes to the network itself. Typical examples are denoising \cite{li_deep_2019}, and normalisation together with regularisation \cite{theate_application_2021}. For exploration, the \(\varepsilon\)-greedy policy is standard when the action set is discrete, and additive Gaussian noise when it is continuous \cite{wu_adaptive_2020, majidi_algorithmic_2022}.

Trading costs are not recorded in the same way across studies. Five of the 18 studies state that their headline result still holds with explicit and implicit costs \cite{huang_financial_2018, gran_deep_2019, yang_deep_2020, jia_lstm-ddpg_2021, zhang_deep_2020}. The others report returns without saying what was deducted, so the reader has to guess whether the figure is gross or net. Training usually uses a train-validation-test split, sometimes with cross-validation. One study trains online and keeps learning as new data arrives \cite{huang_financial_2018}. Of all the studies, this one comes closest to the conditions of live deployment.

The data covers long periods, but at a coarse resolution. Twelve of the 18 studies cover ten years or more, and the median span is eleven years. So a long history is normal in this field. The studies are also alike in resolution, which is low: 17 of the 18 work on bars, 13 of those on daily bars, and only one uses quote-level data \cite{carapuco_reinforcement_2018}. This matters for how costs are counted. A bar reports four prices and hides the order in which the price reached them.

% !TEX root = ../../Main.tex
\section{Research Gap: Positioning This Work}

The eighteen studies agree that a deep reinforcement learning agent can learn a profitable trading policy. They differ, or say nothing, on the conditions in which this holds. The findings above describe what the field does. When they are read against one another, they also show what the field has not done yet. \Cref{Tables:Positioning} sets out the comparison.

% !TEX root = ../Main.tex
\begin{table}[htb!]
\centering
\footnotesize
\caption{How this work compares with the eighteen surveyed articles.}
\label{Tables:Positioning}
\begin{tabularx}{\textwidth}{@{}lXX@{}}
\toprule
\textbf{Dimension} & \textbf{Surveyed literature (18 articles)} & \textbf{This work} \\
\midrule
Market & Stocks 12, Forex 3, ETFs 2, Futures 1 & Forex, the seven major pairs \\
\addlinespace
Data resolution & Bars in 17 of 18, daily in 13; one study uses quotes & Quote-level tick data, with hourly and daily decisions \\
\addlinespace
Data span & Ten years or more in 12 of 18, median eleven years & Eleven years, 2015 to 2026 \\
\addlinespace
Action space & Discrete 15, continuous 3; one sizes the position & Continuous, sets direction and size together \\
\addlinespace
Reward function & Total return 14, Sharpe ratio 3, log return 1 & Log return; models selected by the Calmar ratio and a risk-adjusted composite score \\
\addlinespace
Trading costs & Charged in 5 of 18 & Quoted spread at the moment of the trade plus the commission of a swap-free raw-spread account \\
\addlinespace
Instruments per design & One agent per instrument, or one agent on a basket & One design trained separately on seven pairs, then combined \\
\addlinespace
Reported robustness & The candidate count and the sensitivity to starting conditions are not stated & Five starting balances per pair, with the candidate count stated \\
\bottomrule
\end{tabularx}
\end{table}

The first gap appears only when two of those dimensions are read together. The three studies that trade currencies \cite{carapuco_reinforcement_2018, huang_financial_2018, rundo_deep_2019} all limit the agent to a discrete choice. The three that give the agent a continuous action \cite{majidi_algorithmic_2022, jia_lstm-ddpg_2021, zhang_deep_2020} all trade something else. In the work reviewed here the two properties never appear together. So there is no direct comparison for a currency agent that sets its own position size. This work uses that combination, and for this reason it cannot be measured against an existing result.

The second gap comes from the cost figures. A strategy tested without costs is not tested against the market it claims to trade. In Forex this omission matters more than the reported returns suggest, because the spread is not fixed. It widens when liquidity falls and around releases of economic data. These are the moments when a model trained on an average spread would expect to act. The third gap is about resolution, not length of history. A long sample is already normal in this literature, so it does not set a study apart, but working inside the bar does. A bar reports four prices and hides the order in which the price reached them. That order decides which spread was paid, and whether a stop was reached before a target.

The fourth gap is about evidence, not design. Each study reports the performance of a model it selected. None of the studies reviewed here states how many candidates that model was selected from. None states how the reported figure changes when the evaluation is repeated from a different starting condition. Deep reinforcement learning is sensitive to initialisation. So a single figure cannot separate what comes from the method and what comes from the particular run. A model chosen as the best of many also tells as much about the sample it was chosen from as about the market it traded \cite{bailey_pseudo-mathematics_2014, lopez_de_prado_advances_2018}. The last gap is about how results are combined. The studies reviewed report either one agent trading one instrument, or one agent trading a basket. None of them trains one design separately on each of several instruments and then measures what those agents do when they run together. That is how a trader who holds several currency pairs would deploy them. The individual results do not show what the combination does, because that depends on how the equity curves move against each other.

This work is placed in these gaps. One design is applied without change to all seven major currency pairs. It uses a continuous action that sets the direction and the size of the position together. The evaluation runs on quote-level data over eleven years. It charges the spread quoted at the moment of each trade, plus the commission of a swap-free raw-spread retail account. This is an account that a retail trader can open, and not a simplification made for modelling. Each pair is reported from five different starting balances, not from a single backtest, and the number of candidates the selected model was drawn from is stated. The seven agents are then combined into an equal-weight portfolio. So this work reports both the seven single results and the one result that a trader running all seven would have seen. \Cref{Chapter:ExperimentalResults} gives both.

There is also one point where this work is the same as the literature. The reward is the log return. One of the surveyed studies makes the same choice \cite{majidi_algorithmic_2022}, so this is not a point of difference. The difference is that the reward the agent optimises and the criteria used to select the final model are, on purpose, not the same measure. During training, candidate runs are scored by the Calmar ratio. The model finally published for each pair is the one ranked highest by a composite score. This score is built from percentile ranks within the pair's own pool of candidates. It uses the annualised return, the Sharpe, Sortino, Calmar and Sterling ratios defined in \Cref{Section:PerformanceMeasurement}, the maximum and downside risk, and measures of how the policy behaves. So a policy that earns its return by taking more risk does not win on return alone. \Cref{Chapter:Implementation} sets out both.

The studies reviewed here answer the question of whether a deep reinforcement learning agent can learn a profitable trading policy. They say much less about the conditions in which that answer holds. Each study usually describes the environment its agent was trained in with a sentence or two, and that description is rarely enough to build the environment again.

This work had to build that environment before it could ask the question. An agent trained on a simplified market learns to exploit the simplification. So the value of any result depends on how the environment charges for trading, moves time forward and fills orders. \Cref{Chapter:Implementation} describes that environment and the agent that runs in it.